\section*{A Reader's Map}

\begin{table}[h]
\centering
\footnotesize
\begin{tabularx}{\columnwidth}{>{\RaggedRight\bfseries}p{1.9cm} >{\RaggedRight}X >{\RaggedRight}p{1.2cm}}
\toprule
Part & Reading Purpose & Scope \\
\midrule
I: Decision & Context, motivation, constraints, behavioral psychology & 1 to 12 \\
\addlinespace
II: JTBD & Identify progress sought and forces behind switching & 13 to 21 \\
\addlinespace
III: Acute Pain & Consequential, funded, time-sensitive struggles & 22 to 30 \\
\addlinespace
IV: Segments & Group by decisions rather than decorative personas & 31 to 38 \\
\addlinespace
V: Discovery & Interviews, observation, artifacts, and surveys & 39 to 60 \\
\addlinespace
VI: Preference & Costly commitments, purchases, realized value & 61 to 70 \\
\addlinespace
VII: Buying & Reconstruct competing options and committee votes & 71 to 83 \\
\addlinespace
VIII: Sectors & FMCG, software, consulting, marketplaces, education & 84 to 92 \\
\addlinespace
IX: Bayes Lab & D2C learners, college partners, enterprise buyers & 93 to 117 \\
\addlinespace
X: Decisions & Convert evidence into choices and testable models & 118 to 125 \\
\addlinespace
Appendices & Protocols, registers, and annotated foundations & A to L \\
\bottomrule
\end{tabularx}
\caption{Structural architecture of Chapter 6.}
\label{tab:ch6_readers_map}
\end{table}

\noindent\textbf{How to use this chapter:} Read Parts I to III for the core conceptual model; Parts IV to VI before interviewing or testing commercial offers; Part VII before selling to complex organizations; and Part IX for the continuing Bayes Institute operational case. The appendices provide field protocols intended for repeated use while conducting live customer discovery.

\section*{Prologue: The Customer Who Loved the Idea and Never Bought}

A founder speaks with 30 software engineers. Nearly everyone agrees that artificial intelligence matters immensely. Most affirm that a disciplined, live, mentor-supported technical program would be extraordinarily useful. Several praise the syllabus. Some indicate they would gladly pay a meaningful fee. The founder builds 6 months of instruction, hires rigorous reviewers, designs a polished website, and opens applications. Almost nobody enrolls.

There are many explanations. Perhaps the engineers sincerely valued learning but could not spare 6 months of weekends. Perhaps they already trusted their companies' internal engineering support. Perhaps they feared being judged in an intensive, transparent cohort. Perhaps the actual purchaser would have been their employer. Perhaps their dominant alternative was free personal experimentation with code generation tools, and its benefits were already sufficient. Perhaps the initial respondents were supportive friends being polite. Perhaps there was no immediate, painful reason to act.

None of these explanations can be discovered by asking more people whether AI is important. That question never distinguished \textbf{passive interest from an active purchasing decision}.

Another founder investigates a different way. She asks each engineer about the \textbf{last real occasion} when they had to build something they did not know how to build. She reconstructs the specific task, deadline, consequences, help sought, time invested, money paid, colleagues consulted, alternatives discarded, and outcome. She observes an actual live debugging session. She meets engineers who declined training. She meets a manager who sponsors learning but cannot approve external tools that access proprietary company code. She discovers that some engineers need a 4-hour expert review of a live architectural design, not a course; others need structured accountability because they cannot sustain solitary study; still others need formal credentials the new institution cannot provide. She makes 1 small, bounded, paid diagnostic offer to the subgroup whose actual behavior supports it.

Perhaps this second founder also fails. But she has begun to learn \textbf{which real transaction exists}, rather than which abstract idea sounds attractive.

Chapter 5 taught us that a market is a system of exchanges, not a headline research statistic. Chapter 6 goes inside that system. It asks what causes an individual or an organization to switch from one state to another. The investigation is deliberately rigorous because almost every subsequent decision: product design, differentiation, pricing, distribution, sales, hiring, and capital allocation: depends upon it.

\subsection*{Reading Orientation and Epistemic Rules}

This chapter progresses systematically from decision theory to field research, from individual choices to complex organizational purchases, and then to actionable customer evidence. The Bayes Institute case is an illustration of hypotheses and research design, \textbf{not} a claim that interviews, transactions, educational transformations, institutional contracts, or a fully validated market already exist. All numerical figures are worked pedagogical examples, not survey findings or Indian market benchmarks. Research-derived theories are linked in the references; the canvases, matrices, rules of thumb, and proposed field procedures are managerial syntheses, not scientifically proven recipes. Psychology provides observable tendencies, not universal tricks for manipulating buyers.
